% uid: ""
% school: "lycee"
% level: "1ere"
% chapter: "Théorème de Thalès et sa réciproque"
% title: "Une droite par le centre de gravité"
% difficulty: "Difficile"
% tags: ["thales", "papillon", "milieu", "inegalites"]
% provenance: "own"
% batch: "2026-11-22-revision-1ere-A"

% ---STATEMENT---
$ABC$ est un triangle, $I$ le milieu de $[BC]$ et $G$ le point de $[AI]$ tel que $AG=\dfrac23AI$. Une droite $d$ passant par $G$, mais pas par $A$, coupe $[AB]$ en $M$ et $[AC]$ en $N$. Les parallèles à $d$ passant par $B$ et par $C$ coupent la droite $(AI)$ en $U$ et $V$ ; on admet que $U$ et $V$ sont sur la demi-droite $[AI)$, au-delà de $G$.
% alt: Triangle ABC, milieu I de BC, point G aux deux tiers de AI, droite d passant par G coupant AB en M et AC en N, parallèles à d par B et C coupant AI en U et V
\begin{center}
\begin{tikzpicture}[scale=1.15]
  \coordinate (A) at (3,6); \coordinate (B) at (0,0); \coordinate (C) at (9,0); \coordinate (I) at (4.5,0); \coordinate (G) at (4,2);
  \coordinate (M) at (0.75,1.5); \coordinate (N) at (6.6,2.4); \coordinate (U) at (4.3333,0.6667); \coordinate (V) at (4.6667,-0.6667);
  \draw[thick] (A)--(B)--(C)--cycle; \draw[gray] (A)--(V);
  \draw[thick] ($(M)!-0.35!(N)$)--($(N)!-0.25!(M)$);
  \draw[dashed] (B)--(U); \draw[dashed] (C)--(V);
  \foreach \p/\pos in {A/above,B/below left,C/below right,I/above left,G/above left,M/left,N/right,U/above right,V/below right}{\fill (\p) circle (2pt) node[\pos]{$\p$};}
  \node[axmain] at (-0.55,1.25) {$d$};
\end{tikzpicture}
\end{center}
\begin{enumerate}
  \item Montrer que $I$ est le milieu de $[UV]$, puis que $AU+AV=2\,AI$. [Papillon de sommet $I$ dans les triangles $IBU$ et $ICV$.]
  \item Écrire Thalès de sommet $A$ dans les triangles $ABU$ et $ACV$ pour montrer que $\dfrac{AB}{AM}=\dfrac{AU}{AG}$ et $\dfrac{AC}{AN}=\dfrac{AV}{AG}$.
  \item En déduire que $\dfrac{AB}{AM}+\dfrac{AC}{AN}=3$.
  \item Application : $AB=12$, $AC=10$ et $AM=9$. Calculer $AN$. Que devient la relation lorsque $d\parallel(BC)$ ? lorsque $d=(BG)$ ?
  \item On pose $x=\dfrac{AM}{AB}$, $y=\dfrac{AN}{AC}$ et on note $\mathcal S$ l'aire de $ABC$. Montrer que $\text{aire}(AMN)=xy\,\mathcal S$, puis que $\dfrac49\,\mathcal S\leq\text{aire}(AMN)\leq\dfrac12\,\mathcal S$. Préciser les positions de $d$ pour lesquelles il y a égalité. [Poser $u=\dfrac1x$ ; montrer que $1\leq u\leq2$ et que $xy=\dfrac1{u(3-u)}$.]
\end{enumerate}

% ---SOLUTION---
\begin{prop}[Thalès et papillon]
\textbf{Triangle.} $M\in[AB]$, $G\in[AU]$, $(MG)\parallel(BU)$ : on part de $A$, \ib{$\dfrac{AM}{AB}=\dfrac{AG}{AU}=\dfrac{MG}{BU}$}.

\textbf{Papillon.} Deux droites se coupent en $I$, avec $(BU)\parallel(CV)$ : on part de $I$, \ib{$\dfrac{IB}{IC}=\dfrac{IU}{IV}=\dfrac{BU}{CV}$}.
\end{prop}
\begin{enumerate}
  \item Les droites $(BC)$ et $(UV)=(AI)$ se coupent en $I$. Les droites $(BU)$ et $(CV)$ sont parallèles (toutes deux parallèles à $d$). Papillon de sommet $I$ (triangles $IBU$ et $ICV$) :
  \[\dfrac{IB}{IC}=\dfrac{IU}{IV}=\dfrac{BU}{CV}\]
  $I$ est le milieu de $[BC]$ : $IB=IC$, donc $IU=IV$. Les points $B$ et $C$ sont de part et d'autre de $I$, donc $U$ et $V$ aussi : \ib{$I$ est le milieu de $[UV]$}.

  Les points $U$ et $V$ sont sur la même demi-droite d'origine $A$. Sur cette demi-droite, le milieu de $[UV]$ est à la distance $\dfrac{AU+AV}2$ de $A$ :
  \[AI=\dfrac{AU+AV}2\implies\ib{$AU+AV=2\,AI$}\]

  \item \textbf{Triangle $ABU$.} On a $M\in[AB]$, $G\in[AU]$ (car $U$ est au-delà de $G$) et $(MG)=d\parallel(BU)$. Thalès de sommet $A$ :
  \[\dfrac{AM}{AB}=\dfrac{AG}{AU}\implies\ib{$\dfrac{AB}{AM}=\dfrac{AU}{AG}$}\]
  \textbf{Triangle $ACV$.} On a $N\in[AC]$, $G\in[AV]$ et $(NG)=d\parallel(CV)$. Thalès de sommet $A$ :
  \[\dfrac{AN}{AC}=\dfrac{AG}{AV}\implies\ib{$\dfrac{AC}{AN}=\dfrac{AV}{AG}$}\]

  \item On additionne les deux égalités de la question 2 :
  \[\dfrac{AB}{AM}+\dfrac{AC}{AN}=\dfrac{AU+AV}{AG}=\dfrac{2\,AI}{AG}\]
  Comme $AG=\dfrac23AI$ :
  \[\dfrac{AB}{AM}+\dfrac{AC}{AN}=\dfrac{2\,AI}{\frac23AI}=\ib{$3$}\]

  \item \textbf{Application.} $\dfrac{12}9+\dfrac{10}{AN}=3\implies\dfrac{10}{AN}=3-\dfrac43=\dfrac53\implies AN=\boxed{6}$.

  \textbf{Si $d\parallel(BC)$.} Thalès de sommet $A$ dans $ABI$ (avec $G\in[AI]$) donne $\dfrac{AM}{AB}=\dfrac{AG}{AI}=\dfrac23$ ; de même $\dfrac{AN}{AC}=\dfrac23$. On retrouve $\dfrac32+\dfrac32=3$.

  \textbf{Si $d=(BG)$.} Alors $M=B$ et $\dfrac{AB}{AM}=1$. La relation donne $\dfrac{AC}{AN}=2$ : $N$ est le milieu de $[AC]$ (la droite $(BG)$ est la médiane issue de $B$).

  \item \textbf{Les aires.} Les triangles $AMN$ et $ABN$ ont le même sommet $N$ et leurs bases $[AM]$ et $[AB]$ sont sur la même droite : même hauteur, donc
  \[\dfrac{\text{aire}(AMN)}{\text{aire}(ABN)}=\dfrac{AM}{AB}=x\]
  De même, $ABN$ et $ABC$ ont le même sommet $B$ et leurs bases $[AN]$ et $[AC]$ sont sur la droite $(AC)$ :
  \[\dfrac{\text{aire}(ABN)}{\text{aire}(ABC)}=\dfrac{AN}{AC}=y\]
  \[\text{aire}(AMN)=xy\,\mathcal S\]
  \textbf{La relation.} La question 3 s'écrit $\dfrac1x+\dfrac1y=3$. On pose $u=\dfrac1x$ et $v=\dfrac1y=3-u$.

  Comme $0<x\leq1$ et $0<y\leq1$ : $u\geq1$ et $v\geq1$, donc $3-u\geq1$, c'est-à-dire $1\leq u\leq2$.
  \[xy=\dfrac1{uv}=\dfrac1{u(3-u)}\]
  \textbf{Encadrement de $u(3-u)$.} On écrit
  \[u(3-u)=\dfrac94-\left(u-\dfrac32\right)^2\]
  Pour $1\leq u\leq2$ : $0\leq\left(u-\dfrac32\right)^2\leq\dfrac14$, donc
  \[2\leq u(3-u)\leq\dfrac94\implies\dfrac49\leq xy\leq\dfrac12\]
  \[\ib{$\dfrac49\,\mathcal S\leq\text{aire}(AMN)\leq\dfrac12\,\mathcal S$}\]
  \textbf{Cas d'égalité.} $xy=\dfrac49$ si et seulement si $u=\dfrac32$, c'est-à-dire $x=y=\dfrac23$ : $d$ est parallèle à $(BC)$.

  $xy=\dfrac12$ si et seulement si $\left(u-\dfrac32\right)^2=\dfrac14$, c'est-à-dire $u=1$ ou $u=2$ : $M=B$ et $N$ milieu de $[AC]$ (la médiane $(BG)$), ou bien $N=C$ et $M$ milieu de $[AB]$ (la médiane $(CG)$).
\end{enumerate}
\textbf{Idée :} les deux parallèles à $d$ menées par $B$ et $C$ transportent les rapports $\dfrac{AB}{AM}$ et $\dfrac{AC}{AN}$ sur la droite $(AI)$, où le milieu $I$ les fait s'additionner.

\textbf{Piège :} écrire Thalès dans le triangle $ABC$ : la droite $(MN)$ n'est pas parallèle à $(BC)$ en général.
